% ATTEMPT_STATUS: unresolved
% TOP_PROBLEM: {"problemId": "problem.the-baum-connes-conjecture-for-the-k-theory-of-reduced-group-c-algebras", "canonicalTitle": "The Baum-Connes Conjecture for the K-Theory of Reduced Group C*-Algebras", "releaseRank": 53, "primaryDomain": "algebra_representation_category", "primaryDomainLabel": "Algebra, representation & category theory", "displayStatus": "open", "releaseStatus": "open"}
% !TeX program = lualatex
% Complete problem section copied verbatim from research_batch_51_54.tex.
\documentclass[11pt]{article}
\usepackage[margin=25mm]{geometry}
\usepackage{fontspec}
\setmainfont{Latin Modern Roman}
\usepackage{amsmath,amssymb,mathtools}
\usepackage{unicode-math}
\setmathfont{Latin Modern Math}
\usepackage{enumitem,needspace}
\usepackage{xurl,hyperref}
\hypersetup{colorlinks=true,urlcolor=blue}
\setcounter{secnumdepth}{1}
\setlength{\parindent}{0pt}
\setlength{\parskip}{5pt}
\setlength{\emergencystretch}{3em}
\setlist[itemize]{leftmargin=3em,itemsep=5pt,topsep=4pt,parsep=0pt}
\allowdisplaybreaks
\newcommand{\N}{\mathbb{N}}
\newcommand{\Z}{\mathbb{Z}}
\newcommand{\Q}{\mathbb{Q}}
\newcommand{\R}{\mathbb{R}}
\newcommand{\C}{\mathbb{C}}
\newcommand{\F}{\mathbb{F}}
\newcommand{\A}{\mathbb{A}}
\newcommand{\PP}{\mathbb P}
\newcommand{\E}{\mathbb{E}}
\newcommand{\eps}{\varepsilon}
\newcommand{\dd}{\,\mathrm d}
\newcommand{\sref}[2]{\hyperlink{source:#1:#2}{[S#2]}}
\newcommand{\eref}[2]{\hyperlink{extra:#1:#2}{[E#2]}}
% Turn the next switch off for a shorter reading copy. The quotations preserve
% every exactTarget field, including qualifications omitted from a short summary.
\newif\ifshowcatalogue
\showcataloguetrue
\newcommand{\cataloguescope}[1]{\ifshowcatalogue
 \paragraph{Exact scope in the supplied catalogue.}
 \begin{quote}\small #1\end{quote}\fi}
\begin{document}
\setcounter{section}{52}
% ================================================================
% problemId: problem.the-baum-connes-conjecture-for-the-k-theory-of-reduced-group-c-algebras
% source releaseRank: 53; source releaseStatus: open
% exactTarget SHA-256: 1e80ca06bf1a687d4292bebe3976cb6de48e445e8334000523e6b6ad9407b631
\clearpage
\section{\texorpdfstring{The Baum-Connes Conjecture for the K-Theory of Reduced Group C*-Algebras}{The Baum-Connes Conjecture for the K-Theory of Reduced Group C*-Algebras}}\label{53}
\subsection{Definitions and mathematical statement}
Let $\underline EG$ be a classifying space for proper $G$-actions: fixed-point sets for compact subgroups are contractible and those for noncompact subgroups are empty. Topological $K$-theory here is the compact-support/direct-limit equivariant $K$-homology of $\underline EG$. The reduced group $C^*$-algebra $C_r^*(G)$ is the operator-norm completion of convolution operators in the left regular representation. The assembly map is
\[
 \mu_G:K_*^{\rm top}(G)\longrightarrow K_*(C_r^*(G)),\qquad *=0,1.
\]
The conjecture says it is an isomorphism for every second-countable locally compact $G$. This is the coefficient-free target; adding arbitrary coefficient algebras changes the conjecture.
\par\smallskip\textit{Formulation and concept references:} \sref{53}{1}.
\cataloguescope{Prove the Baum-Connes conjecture for every second-countable locally compact group G: the assembly map K\_top(G) \ensuremath{\to} K(C*\_r(G)) from the topological K-theory of the group, built from the classifying space for proper actions, to the K-theory of the reduced group C*-algebra is an isomorphism.}
\subsection{Short English statement}
Does geometric equivariant $K$-homology compute the operator-algebraic $K$-theory of every group in the stated class?
% BEGIN RESEARCH ADDITION ATTEMPT-53
\subsection{Research attempt}

\paragraph{Current frontier and source audit.}
The non-discrete formulation was checked against Section~4.1.6 of
L\"uck--Reich. Their Sections~5.1.1--5.1.2 record the a-T-menable and
almost-connected cases, among other positive results. The target here
still quantifies over every second-countable locally compact group and
uses only the coefficient algebra $\C$. \sref{53}{1}
Nishikawa's Theorem~A and Corollary~13.4 provide a controlled-algebra
reformulation for this full locally compact scope, not just for discrete
groups. \eref{53}{1}

\textbf{Editorial status note.} The paper listed as [S4] was withdrawn
in arXiv version~2 on 12 May 2026. Its author reports a missing norm-control
argument in the $K$-theory component. Its advertised conclusion is
therefore not used as an established theorem here. \sref{53}{4},
\eref{53}{2}
Meyer's 2025 note concerns the hypotheses of an extension permanence
theorem with coefficients. It disproves an attempted simplification of
those hypotheses, not the coefficient-free conjecture in this section.
\sref{53}{5}, \eref{53}{3}

\paragraph{Attack route.}
Three approaches lead to different missing statements. In the
Dirac--dual-Dirac approach, an appropriate gamma element gives a split
injection, but surjectivity still requires that its descended action be
the identity on $K_*(C_r^*(G))$. Requiring a gamma element equal to the
identity in equivariant $KK$-theory for every group would be stronger:
it would give arbitrary coefficients as well. One cannot silently add
that stronger requirement, in view of coefficient counterexamples.
\eref{53}{1}, \eref{53}{3}

An induction using group extensions would need the full finite-subgroup
preimage conditions in the appropriate permanence theorem; the shortcut
of testing only a normal subgroup is specifically unsafe. \eref{53}{3}
The route tested below is instead to regularize representatives of
$K$-classes until geometric control becomes small, then recover exact
projections or unitaries by functional calculus. Its missing ingredient
is quantitative: the spectral gap needed for the recovery must not
collapse during regularization. This is tested by an exact example, not
by an assumption that norm-small coefficients imply controlled
$K$-theory.

\paragraph{Attempt: identify the obstruction algebra in the full scope.}
Here is the established bridge used, with coefficients specialized to
$\C$. For a proper $G$-space $X$ and a universal $X$-$G$ Hilbert module
$H_X$ in Nishikawa's Definition~5.2, $RL_c^*(H_X)$ is the completion of
bounded, norm-continuous, $G$-continuous compact-operator-valued
paths $T_t$, $t\ge1$,
with uniformly compact support and
$\|[T_t,\phi]\|\to0$ for $\phi\in C_0(X)$.
Let $RL_c^0(H_X)$ be its evaluation-at-one kernel. Set
\begin{equation}\label{ra:53:kernel}
 I_G=RL_c^0(H_{\underline EG})\rtimes_r G.
\end{equation}
For non-locally-compact models of $\underline EG$, use the source's
direct-limit convention over $G$-compact proper subspaces and compatible
modules. Corollary~13.4 identifies the required condition as
\begin{equation}\label{ra:53:obstruction}
 \mu_G\text{ is an isomorphism}
 \quad\Longleftrightarrow\quad K_0(I_G)=K_1(I_G)=0.
\end{equation}
This is a known reformulation, not a newly proved reduction. The
semisplit evaluation sequence is what permits the reduced-crossed-product
argument; exactness of reduced crossed products for arbitrary sequences
is not being assumed. \eref{53}{1}

A concrete sufficient analytic programme is now visible. For every
projection $p$ in a matrix unitization of $I_G$, with scalar part $e$,
construct a norm-continuous path $a_s$, $s\in[0,\infty)$, in the same
matrix unitization, satisfying
\begin{equation}\label{ra:53:program0}
 \begin{gathered}
 a_0=p,\qquad a_s=a_s^*,\qquad a_s-e\in M_N(I_G),\qquad
 \lim_{s\to\infty}\|a_s-e\|=0,\\
 \sup_{s\ge0}\|a_s^2-a_s\|\le\varepsilon_0<\tfrac14.
 \end{gathered}
\end{equation}
Stabilization is allowed. For every unitary $u$ with scalar part one,
seek norm-continuous paths $b_s$, starting at $u$, tending in norm
to one, and satisfying
\begin{equation}\label{ra:53:program1}
 b_s-1\in M_N(I_G),\qquad
 \sup_s\max\{\|b_s^*b_s-1\|,\|b_sb_s^*-1\|\}
 \le\varepsilon_1<1.
\end{equation}
The $s$ here is a homotopy variable. It is different from the localization
variable $t$ hidden in $I_G$; controlling each separately is not enough
unless the resulting paths are continuous in the algebra norm.

\paragraph{Attempt: prove the functional-calculus step.}
\textbf{Lemma 53.1 (uniform spectral recovery).} In a unital
$C^*$-algebra, a self-adjoint element $a$ with
$\|a^2-a\|\le\varepsilon<1/4$ has spectrum separated from $1/2$ by at
least $\sqrt{1-4\varepsilon}/2$. Its spectral projection
$p(a)=\mathbf1_{(1/2,\infty)}(a)$ is well defined by continuous
functional calculus on its spectrum and
\begin{equation}\label{ra:53:recovery}
 \|p(a)-a\|\le
 r_\varepsilon:=\frac{1-\sqrt{1-4\varepsilon}}2.
\end{equation}
The assignment is norm-continuous on any norm-bounded family with the
same error bound. It respects a quotient in which $a$ maps to a projection.

\textit{Proof.} The spectral theorem implies
$|\lambda^2-\lambda|\le\varepsilon$ for every spectral value $\lambda$.
Inside $[0,1]$ this excludes the interval
$(r_\varepsilon,1-r_\varepsilon)$; the possible excursions below zero
or above one have size
$(\sqrt{1+4\varepsilon}-1)/2\le r_\varepsilon$.
This proves the spectral gap and the distance bound. Choose a single
continuous function equal to zero on the lower allowed band and one on
the upper band. Its functional calculus gives $p(a)$ for every element
of the family. Uniform polynomial approximation of this function on a
common compact interval proves norm-continuity. Functoriality under a
quotient, and the values zero and one at those scalars, prove the last
assertion. \hfill$\square$

The analogous recovery of a unitary is
\[
 v(b)=b(b^*b)^{-1/2}.
\]
Both inequalities in \eqref{ra:53:program1} ensure invertibility; the
uniform lower bound $b^*b\ge(1-\varepsilon_1)1$ makes inverse square roots
continuous with uniform control. These statements follow from scalar
continuous functional calculus and do not require an assembly theorem.

Applying Lemma~53.1 to \eqref{ra:53:program0} would turn $a_s$ into a
projection homotopy from $p$ to $e$, with the endpoint at infinity supplied
by norm convergence. Applying the polar construction to
\eqref{ra:53:program1} would give a unitary homotopy from $u$ to one.
Thus the proposed programme would kill both $K$-groups in
\eqref{ra:53:obstruction}. This implication is proved; the paths in
\eqref{ra:53:program0}--\eqref{ra:53:program1} have not been constructed.

In a discrete equivariant Roe/localization model there is a useful
additional check. Products of $k$ finite-propagation operators have
propagation at most the sum of their propagations. Hence a fixed-degree
polynomial in a path whose propagation tends to zero again has that
property. The uniform polynomial approximation in the proof above
therefore keeps spectral recovery inside the \emph{completed}
localization algebra. It does not show that an exact recovered projection
itself has literal finite propagation. Confusing the dense controlled
subalgebra with its norm completion would be an error.

\paragraph{Attempt: test a natural damping regularization.}
Consider already $G=\Z$, so $C_r^*(G)\cong C(S^1)$, and write $u(z)=z$.
For $s\ge0$, the Poisson averaging operator satisfies
\[
 M_s(u^k)=e^{-s|k|}u^k.
\]
It is unital and completely positive: for $s>0$ it is averaging of
rotations against the nonnegative Poisson kernel of radius $e^{-s}$,
whose integral is one, and the same averaging works entrywise on
matrices. The identity case is $s=0$. It therefore looks attractive as a
regularizer that suppresses nontrivial group coefficients.

The following calculation breaks that proposal's spectral-control step.
In $M_2(C(S^1))$ put
\[
 p=\frac12\begin{pmatrix}1&u\\u^*&1\end{pmatrix},\qquad
 a_s=M_s^{(2)}(p)=\frac12
 \begin{pmatrix}1&e^{-s}u\\e^{-s}u^*&1\end{pmatrix}.
\]
Direct multiplication, using $u^*u=uu^*=1$, gives
\begin{equation}\label{ra:53:damping}
 p^2=p,\qquad
 a_s^2-a_s=-\frac{1-e^{-2s}}4 I_2,
 \qquad
 \operatorname{spec}(a_s)=
 \left\{\frac{1-e^{-s}}2,\frac{1+e^{-s}}2\right\}.
\end{equation}
For each finite $s$ the idempotence error is strictly below $1/4$.
Nevertheless its supremum is $1/4$, and the spectral gap is $e^{-s}$,
which tends to zero. Functional calculus does not retain the coefficient
suppression:
\begin{equation}\label{ra:53:undamping}
 \mathbf1_{(1/2,\infty)}(a_s)=p\qquad(s<\infty).
\end{equation}
The exact spectral projection restores the original off-diagonal $u$.
Moreover $a_s\to I_2/2$, which is not a projection. For $K_1$ the same
phenomenon is even simpler: $M_s(u)=e^{-s}u$ is invertible at every finite
$s$, but its polar unitary is always $u$, while the inverse norm grows
as $e^s$.

This is not a counterexample to Baum--Connes, nor even a nonzero
obstruction class. Indeed the displayed projection is explicitly
homotopic to the constant rank-one projection through
\[
 p_\theta=
 \begin{pmatrix}
 \cos^2\theta&\cos\theta\sin\theta\,u\\
 \cos\theta\sin\theta\,u^*&\sin^2\theta
 \end{pmatrix},\qquad 0\le\theta\le\pi/4.
\]
Every $p_\theta$ is a projection, $p_0=\operatorname{diag}(1,0)$ and
$p_{\pi/4}=p$. The example rejects the specific coefficient-damping
argument, even in a group for which the conjecture is established; it
does not reject the existence of a different geometric homotopy.

\paragraph{Stress test.}
There are three distinct uniformities. Spectral separation must be
uniform along the proposed homotopy. Localization estimates must hold
in the completed algebra, rather than separately for each $t$. Finally,
a direct-limit argument may enlarge the proper subspace for an individual
class, but cannot let that subspace run off to infinity during a purported
homotopy without proving that the homotopy represents a class in the
required direct limit. None is supplied by pointwise decay of matrix
coefficients.

The exact example \eqref{ra:53:damping} is a counterexample to the
intermediate claim that pointwise almost-idempotence below $1/4$ is
sufficient for uniform spectral recovery during regularization. Its
failure is independent of numerical accuracy. The withdrawn norm-control
argument in [S4] is a separate source warning, not evidence that the
calculation here reproduces that paper's particular missing step.
\eref{53}{2}

A proof only lifting $K$-classes to an assembly image would address
surjectivity, not injectivity. The obstruction-algebra formulation and
both proposed contractions were included to avoid that omission.
Likewise, calculations with $\Z$ test a mechanism but do not reduce all
locally compact groups to discrete ones. The general formulation remains
\eqref{ra:53:kernel}--\eqref{ra:53:obstruction}.

\paragraph{Outcome.}
\textbf{Outcome: Exploratory approach with unresolved gap.}

The functional-calculus recovery lemma and the explicit damping failure
were proved. Together they identify a precise defect in one possible
controlled-algebra attack: spectral flattening can erase the very gap
needed to pass back from approximate to exact $K$-theory representatives.
The known localization equivalence is used as a framework, not presented
as new progress.

The unresolved step is the construction, for every group in the original
scope, of geometric homotopies with genuinely uniform norm and spectral
control, or another method that kills the same obstruction classes.
No such construction, universal assembly inverse, or coefficient-free
counterexample was obtained. In particular, neither an extension claim
with altered hypotheses nor a withdrawn quantitative theorem is used to
fill that gap.
% END RESEARCH ADDITION ATTEMPT-53
\subsection{Sources}
\begin{itemize}\raggedright
\item[\textnormal{[S1]}]\hypertarget{source:53:1}{} W. Lück, H. Reich, in Handbook of K-theory, Springer, 2005.
\url{https://arxiv.org/abs/math/0402405}.
{\small\textit{Catalogue formulation source.}}
\item[\textnormal{[S2]}]\hypertarget{source:53:2}{} There are no known counter-examples to the Baum-Connes conjecture.
\url{https://www.fields.utoronto.ca/talks/There-are-no-known-counter-examples-to-Baum-Connes-conjecture}.
{\small\textit{Catalogue research/status reference; not a proof endorsement.}}
\item[\textnormal{[S3]}]\hypertarget{source:53:3}{} Expanders and K-theory for group C* algebras.
\url{https://www.fields.utoronto.ca/talk-media/2/78/29/slides.pdf}.
{\small\textit{Catalogue research/status reference; not a proof endorsement.}}
\item[\textnormal{[S4]}]\hypertarget{source:53:4}{} Controlled Analytic Properties and the Quantitative Baum-Connes Conjecture.
\url{https://arxiv.org/abs/1908.02131}.
{\small\textit{Catalogue research/status reference; not a proof endorsement.}}
\item[\textnormal{[S5]}]\hypertarget{source:53:5}{} The Baum-Connes conjecture for extensions.
\url{https://arxiv.org/abs/2508.05726}.
{\small\textit{Catalogue research/status reference; not a proof endorsement.}}
% BEGIN RESEARCH ADDITION REFERENCES-53
\item[\textnormal{[E1]}]\hypertarget{extra:53:1}{} Shintaro Nishikawa, \emph{Crossed product approach to equivariant localization algebras}, arXiv:2108.11235v2, Theorem A, Definition 5.2 and Corollary 13.4; see also the introduction on gamma elements.
\url{https://arxiv.org/abs/2108.11235v2}.
{\small\textit{Added research source: Localization/crossed-product formulation for second-countable locally compact groups, including the kernel criterion and the semisplit evaluation sequence. Consulted 24 September 2026.}}
\item[\textnormal{[E2]}]\hypertarget{extra:53:2}{} Martin Finn-Sell, \emph{Controlled Analytic Properties and the Quantitative Baum--Connes Conjecture}, arXiv:1908.02131v2, author withdrawal notice dated 12 May 2026.
\url{https://arxiv.org/abs/1908.02131v2}.
{\small\textit{Added research source: Withdrawn: the notice reports a missing norm-control argument in the K-theory component. Cited for source status only, not as an established quantitative theorem. Consulted 24 September 2026.}}
\item[\textnormal{[E3]}]\hypertarget{extra:53:3}{} Ralf Meyer, \emph{The Baum--Connes conjecture for extensions}, arXiv:2508.05726v2, 29 August 2025, two-page note.
\url{https://arxiv.org/abs/2508.05726v2}.
{\small\textit{Added research source: Counterexample to an overstrong extension-permanence assertion with coefficients. This is not a counterexample to the coefficient-free target. Consulted 24 September 2026.}}
% END RESEARCH ADDITION REFERENCES-53
\end{itemize}

\end{document}
